
\subsection{电容与偏转}

\begin{tabularx}{\columnwidth}{XX}
  电容定义&$C=\dfrac{Q}{U}$\\
  平行板电容器&$C=\varepsilon_0\varepsilon_r\dfrac{S}{d}$\\
  电容器能量&$E_C=\dfrac12CU^2=\dfrac{Q^2}{2C}=\dfrac12QU$\\
  加速电场&$qU=\dfrac12mv^2$\\
  偏转电场加速度&$a=\dfrac{qE}{m}=\dfrac{qU}{md}$\\
  偏转位移&$y=\dfrac12at^2$\\
\end{tabularx}

\begin{formulanotebox}
平行板电容器动态分析先判断约束：接电源时 $U$ 不变，断开电源时 $Q$ 不变，再由 $C=\varepsilon_0\varepsilon_r\dfrac{S}{d}$ 推出其他量。

带电粒子在偏转电场中通常分解为沿初速度方向的匀速直线运动和垂直电场方向的匀加速运动。
\end{formulanotebox}

\begin{keypointbox}[title={\strut\textcolor{white}{\bfseries 重点知识点：电容器动态分析}}]
\textbf{接电源}：电压由电源维持不变，改变板间距、正对面积或介质时，先判断 $C$，再判断 $Q=CU$ 和 $E=\frac{U}{d}$。\\

\textbf{断开电源}：电荷量不变，改变结构时，先判断 $C$，再判断 $U=\frac{Q}{C}$。平行板中若只改变板间距，$E=\frac{Q}{\varepsilon_0\varepsilon_r S}$ 常保持不变。\\

\textbf{充放电}：充电初始电流较大，随后逐渐减小；稳定后电容器支路可看成断路。
\end{keypointbox}

\begin{casetypebox}[title={\strut\textcolor{white}{\bfseries 常见题型：电场加速与偏转}}]
\begin{itemize}[leftmargin=1.5em]
  \item 加速区优先用动能定理：$qU=\Delta E_k$。
  \item 偏转区先求运动时间 $t=\frac{L}{v_0}$，再由 $y=\frac12\frac{qU}{md}t^2$ 求偏移。
  \item 离开极板后的粒子做匀速直线运动，速度方向由出射瞬间的合速度决定。
  \item 固定水平位移时，若需要最小合速度，常利用水平、竖直分速度关系判断临界角。
\end{itemize}
\end{casetypebox}
